\section{Homology}

Homology is a hot topic when talking about polyhedra.
Here is a quick account.

\subsection{simplicial homology}
We will define the homology of simplicial complexes. 
A full account can be found in any algebraic topology book, such as \cite{at1}, \cite{at2}.

Providing proofs of these basic facts is beyond the scope of this article, thus we will avoid them completely and 
refer to the forementioned books for any proof of theorems stated below. 

From now on when we call a simplicial complex $(V,\Sigma)$ we also fix an order on the vertices, and build simplices respecting that order.

\begin{definition}[Simplicial $k-$chain]\hfill 

Let $(V,\Sigma)$ be a simplicial complex. We call $\Sigma_k=\set{\sigma\in \Sigma:\#(\sigma)=k}$

We call $C_k$ the \textbf{simplicial k-chain} group 
\[\bigoplus_{\sigma\in \Sigma_k}\Z [\sigma],\]
i.e. the free abelian groups generated by the $k-$simplices of $\Sigma$.
\end{definition}

\begin{definition}[Boundary Map]\hfill

Let $(V,\Sigma)$ be a simplicial complex, $\Sigma\ni\sigma=(v_1,...,v_k)$ be a $k-simplex$.

We define the \textbf{boundary map} 
\[\partial_k:C_k\to C_{k-1},\ \partial_k(\sigma)=\sum_{i=1}^k(-1)^i(v_1,...,\hat{v_i},...,v_k)\]

We define two subgroups of $C_k$
\[Z_k:=\ker(\partial_k), B_k:=\im(\partial_{k+1})\]
\end{definition}

\begin{definition}[Homology groups]\hfill

We define the \textbf{k$^{\text{th}}$ homology group}
\[H_k(\Sigma)=Z_k/B_k\]
\end{definition}
\begin{definition}[Homology chain]
    
Together with the boundary maps, we have a long exact sequence 
\[H_0(\Sigma)\to H_1(\Sigma)\to H_2(\Sigma)\to\cdots\]

That we call such a sequence $H_{\bullet}(\Sigma)$ the \textbf{homology chain} of $\Sigma$.

We say that a morphism of homology chains $f_\bullet:H_\bullet(\Sigma)\to H_\bullet(\Sigma')$ is a collection of
group homomorphisms $f_i:H_k(\Sigma)\to H_k(\Sigma')$ such that every square
\[\begin{tikzcd}[ampersand replacement=\&]
	{H_0(\Sigma)} \& {H_1(\Sigma)} \& {H_1(\Sigma')} \& \cdots \\
	{H_0(\Sigma')} \& {H_1(\Sigma')} \& {H_2(\Sigma')} \& \cdots
	\arrow[from=1-1, to=1-2]
	\arrow["{f_0}"', from=1-1, to=2-1]
	\arrow[from=1-2, to=1-3]
	\arrow["{f_1}"', from=1-2, to=2-2]
	\arrow[from=1-3, to=1-4]
	\arrow["{f_2}"', from=1-3, to=2-3]
	\arrow[from=2-1, to=2-2]
	\arrow[from=2-2, to=2-3]
	\arrow[from=2-3, to=2-4]
\end{tikzcd}\]
commutes.
\end{definition}

The intuitive idea is that the rank of $H_k$ counts \textit{how many $k-$dimensional holes} our complex has, with $H_0$ counting connected components.

\begin{example}\hfill 

    \begin{itemize}
        \item The empty triangle has $H_0\simeq \Z$, $H_1\simeq \Z$ and any subsequent $H_k\simeq 0$
        \item The filled triangle has $H_0\simeq \Z$ and all other $H_k\simeq 0$.
        \item The empty tetrahedron has $H_2\simeq \Z$
        \item $\vdots$
    \end{itemize}
\end{example}

%insertpicturehere
\subsection{Modally undefinable homology groups}
While there's little hope to give much information about homology through a modal language we shall note that our constructions map neatly
onto the homology chain. 
In particular 
\begin{proposition}
$H_\bullet(\Sigma\sqcup\Sigma')\simeq H_\bullet(\Sigma)\oplus H_\bullet(\Sigma')$
\end{proposition}

\begin{proposition}
If $\Sigma\leq\Sigma'$ is a closed subcomplex, then there's a map $H_\bullet(\Sigma)\to H_\bullet(\Sigma')$
\end{proposition}

\begin{proposition}
If $f:\Sigma\to\Sigma'$ is a simplicial map, then there's a bisimplicial map $f_*:H_\bullet(\Sigma)\to H_\bullet(\Sigma')$
\end{proposition}

This means
\begin{lemma}
The homologies of a modally definable class of polyhedra must be closed under direct sum and images.

Or, in the more useful direction:

If the homology groups of a class of polyhedra is \textbf{not} closed under direct sums and images, then the class is not modally definable.
\end{lemma}
\begin{proof}
Follows from the three propositions.
\end{proof}

\begin{corollary}
The class of genus $n$ polyhedra is not modally for any positive $n$ 
\end{corollary}
\begin{proof}
    $H_1(\Sigma)\simeq\Z^n$ for all $\Sigma\in C$. 
    As $\Z^n\oplus \Z^n\neq \Z^n$, the class is not closed under direct sums, thus not modally definable.
\end{proof}